\section*{Author's foreword}
{\small\sffamily\color{copper}FIRST-PERSON DRAFT FOR ARYAN AYYAR'S REVIEW}\par\medskip
Long before I had a research question, I had an image: a large aircraft turning above Hong Kong, close enough to the city for the buildings to seem part of the approach. I was fascinated by aircraft as a child, particularly by the engineering that allowed something so large to move with such apparent precision. Kai Tak gave that fascination a setting. The Checkerboard approach looked like a problem that should not have had enough room for its solution.

Watching Richard Hammond's \emph{Engineering Connections} encouraged a particular way of looking at such things. The spectacle was an invitation to ask how the pieces worked together. A remarkable structure or machine became more interesting when its mechanisms were explained. I wanted the same explanation for Kai Tak. How much of the turn came from the aircraft's capabilities, how much from the shape of the approach, and how much depended on the person flying it?

The imagery and music of \emph{Enter the Dragon} belong to that early fascination too. In my memory, aircraft and historical Hong Kong became associated with a city whose geography seemed unusually present in its everyday life. That association is personal, not documentary evidence about a particular flight or filming location. It helped make the question memorable; it could not answer it.

Years later, while studying for a Master of Finance at Hong Kong Polytechnic University and living in Tsim Sha Tsui, I became more attentive to the city's geography and engineering history. The harbour, the hills and the density of the built environment made the old airport feel less like an isolated aviation curiosity. It was part of a larger question about how Hong Kong made demanding arrangements work. My own time in the city came after Kai Tak had closed. Arriving at the modern airport at Chek Lap Kok was not experiencing the former Runway 13 approach.

I returned to the question with tools developed through research in financial economics and mathematics. I am not a trained aerospace engineer. That matters: an equation I can solve is not automatically an adequate account of an aircraft, and a convincing reconstruction is not automatically a survey. The attraction of the project lies partly in discovering where the mathematics is useful and where it must wait for better evidence.

This paper follows that process. It begins with the airport and moves towards a complete mathematical landing, including the less photogenic business of stopping. Along the way, it asks what can be recovered from dated photographs, maps and published flight records, and what those sources leave unresolved. I hope it preserves something of the original curiosity without using wonder as an excuse for imprecision. The approach becomes more impressive, not less, when we can explain some of its demands and say honestly which questions we have not yet answered.

\vfill{\small This foreword records the author's supplied motivation. It is not a historical source or a claim of professional aerospace qualification.}
\clearpage
